\uuid{szCR}
\titre{Équation homogène/solution particulière}
\auteur{hanani ; mzari}
\organisation{exo7}
\licence{CC BY-NC-SA 4.0}
\datecreate{2020-01-20}
\niveau{L1}
\module{Analyse}
\chapitre{Equation différentielle}
\sousChapitre{Résolution d'équation différentielle du premier ordre}
\difficulte{3}
\theme{Equations du premier ordre, équation homogène, solution particulière}
\exo7qcmid{lille-2.equations-differentielles.011}

\contenu{
\begin{question}
\qclassification{225.01}
\qtags{motcle=équation homogène/solution particulière}
On considère l'équation différentielle $(E)$ : $\displaystyle (1+x^2)y'+2xy=2x$. Parmi les affirmations suivantes cocher celles qui sont vraies :
\begin{answers}
\good{La solution générale de l'équation homogène est : $\displaystyle y=\frac{k}{1+x^2}$, $k\in \Rr$.}
\good{La solution de $(E)$ vérifiant $y(0)=0$ est $\displaystyle y=\frac{x^2}{1+x^2}$.}
\bad{$(E)$ n'admet pas de solution vérifiant $y'(0)=0$.}
\bad{$(E)$ admet une unique solution vérifiant $y'(0)=0$.}
\end{answers}
\begin{explanations}
La solution générale de l'équation homogène est : $\displaystyle Y=\frac{k}{1+x^2}$, $k\in \Rr$, et $y_0=1$ est une solution particulière. La solution générale de $(E)$ est :
$$y=1+\frac{k}{1+x^2},\quad k\in \Rr.$$
La condition $y(0)=0$ donne $k=-1$. Par ailleurs, en posant $x=0$ dans $(E)$, on obtient $y'(0)=0$. Donc toute solution de $(E)$ vérifient la condition $y'(0)=0$.
\end{explanations}
\end{question}
}
